\scanpage{230}
\begin{center}- \mspage{191} -\end{center}\margin{IX}
\uline{Remarques} 14.1.12.\quad a) On notera que l'unicité de $e$ rendant commutatif
le diagramme précédent est immédiate, l'espace rigide-analytique $G_K^{\mathrm{an}}$
étant connexe. L'existence de $e$, par contre, n'est nullement triviale,
et nous nous contentons à ce sujet de renvoyer à la note citée de
M. RAYNAUD. D'autre part, la définition donnée dans loc. cit. de
l'homomorphisme $i^{\mathrm{an}}$ est également de nature transcendante, bien qu'il
puisse être interprété comme étant un objet de nature essentiellement
algébrique (14.1.5) (une fois la construction de $G$ lui-même, donc de
$G_K$, acquise (ce qui se fait par les constructions de \uline{géométrie formelle}
du par. 7)). Je n'en connais aucune description purement algébrico-géométrique,
mais proposerai plus bas (14.4) une caractérisation
``arithmétique'', qu'on obtiendra en se ramenant à une situation de type
fini sur l'anneau $\mathbb Z$. Il semble donc que l'homomorphisme (14.1.5) soit,
du point de vue de la géométrie algébrique, un élément de structure
plus caché que ceux rencontrés au cours des paragraphes précédents.

b) Comme expliqué plus haut, l'homomorphisme $i$ est canoniquement
associé à une trivialisation bien déterminée de la biextension $E_K$ de
$(\underline M_K,\underline M'_K)$, ou encore à un homomorphisme bien déterminé (14.1.6) relevant $p'_K$.
On montre alors que ce dernier n'est autre que l'homomorphisme qui
intervient dans la suite exacte
\[
(14.1.13)\qquad0\longrightarrow\underline M'_K{}^{\mathrm{an}}
\xrightarrow{i'{}^{\mathrm{an}}}G'_K{}^{\mathrm{an}}\longrightarrow A'_K{}^{\mathrm{an}}\longrightarrow0
\]
analogue à (14.1.7), associée au schéma abélien dual $A'_K$ de $A_K$ (à moins
que ce ne soit l'opposé ???).
\margin{509}
\margin{509}
\margin{229}
